\documentclass[a4paper,12pt]{exam}
\usepackage[T1]{fontenc}
\usepackage{lmodern}
\usepackage{comment}
\usepackage{fancyvrb}
\usepackage{graphicx}
\usepackage{geometry}
\usepackage{amssymb,amsmath}
\usepackage{ifxetex,ifluatex}
\usepackage{fixltx2e} % provides \textsubscript
% use upquote if available, for straight quotes in verbatim environments
\IfFileExists{upquote.sty}{\usepackage{upquote}}{}
\ifnum 0\ifxetex 1\fi\ifluatex 1\fi=0 % if pdftex
  \usepackage[utf8]{inputenc}
\else % if luatex or xelatex
  \ifxetex
    \usepackage{mathspec}
    \usepackage{xltxtra,xunicode}
  \else
    \usepackage{fontspec}
  \fi
  \defaultfontfeatures{Mapping=tex-text,Scale=MatchLowercase}
  \newcommand{\euro}{€}
\fi
% use microtype if available
\IfFileExists{microtype.sty}{\usepackage{microtype}}{}
\usepackage{longtable,booktabs}
\usepackage{graphicx}
% Redefine \includegraphics so that, unless explicit options are
% given, the image width will not exceed the width of the page.
% Images get their normal width if they fit onto the page, but
% are scaled down if they would overflow the margins.
\makeatletter
\def\ScaleIfNeeded{%
  \ifdim\Gin@nat@width>\linewidth
    \linewidth
  \else
    \Gin@nat@width
  \fi
}
\makeatother
\let\Oldincludegraphics\includegraphics
{%
 \catcode`\@=11\relax%
 \gdef\includegraphics{\@ifnextchar[{\Oldincludegraphics}{\Oldincludegraphics[width=\ScaleIfNeeded]}}%
}%
\ifxetex
  \usepackage[setpagesize=false, % page size defined by xetex
              unicode=false, % unicode breaks when used with xetex
              xetex]{hyperref}
\else
  \usepackage[unicode=true]{hyperref}
\fi
\hypersetup{breaklinks=true,
            bookmarks=true,
            pdfauthor={},
            pdftitle={Worksheet: Streams},
            colorlinks=true,
            citecolor=blue,
            urlcolor=blue,
            linkcolor=magenta,
            pdfborder={0 0 0}}
\urlstyle{same}  % don't use monospace font for urls
\setlength{\parindent}{0pt}
\setlength{\parskip}{6pt plus 2pt minus 1pt}
\setlength{\emergencystretch}{3em}  % prevent overfull lines
%\setcounter{secnumdepth}{0}

\geometry{
  %includeheadfoot,
  margin=2.54cm
}

\title{Exercises --- Week Seven}
\date{Spring term 2016}
\author{Optional exercise --- Functional Sets}

\begin{document}
\maketitle

You need to have installed the current Scala distribution and the current Java distribution
before commencing these exercises.
Don't forget the documentation for both distributions as they will come in very useful.

Any source code examples are available under the \verb!week07/funsets! folder.

The exercise is based upon one from one of the EPFL courses on functional programming techniques, 
later reused in Martin Odersky's \url{coursera.org} course.

\section*{Preamble}

In this exercise, you will work with a functional representation of sets based on the 
mathematical notion of characteristic functions. 
The goal is to gain practice with higher-order functions.
You should write your solutions by completing the stubs in the \verb!FunSets.scala! file.
You will work with sets of integers.

As an example to motivate our representation, how would you represent the set of all 
negative integers? You cannot list them all\ldots one way would be so say: ``
if you give me an integer, I can tell you whether it’s in the set or not: for 3, 
I say ‘no’; for -1, I say yes''.

Mathematically, we call the function which takes an integer as argument and which 
returns a boolean indicating whether the given integer belongs to a set, 
the \emph{characteristic} function of the set. 
For example, we can characterise the set of negative integers by the characteristic 
function \verb!(x: Int) => x < 0!.

Therefore, we choose to represent a set by its characteristic function and 
define a type alias for this representation:
\begin{Verbatim}
type Set = Int => Boolean	
\end{Verbatim}
Using this representation, we define a function that tests for the presence of a value in a set:
\begin{Verbatim}
def contains(s: Set, elem: Int): Boolean = s(elem)	
\end{Verbatim}

Write your own tests! For this assignment, we don’t give you tests but 
instead the \verb!FunSetSuite.scala! file contains hints on how to write your own 
tests for the assignment.

\begin{questions}
\question
Let’s start by implementing basic functions on sets.
\begin{parts}
\part
Define a function which creates a singleton set from one integer value: 
the set represents the set of the one given element. Its signature is as follows:	
\begin{Verbatim}
def singletonSet(elem: Int): Set	
\end{Verbatim}
Now that we have a way to create singleton sets, we want to define a function that 
allow us to build bigger sets from smaller ones.
\part
Define the functions \verb!union!, \verb!intersect!, and \verb!diff!, 
which take two sets, and return, respectively, their union, intersection and differences. 
\verb!diff(s, t)! returns a set which contains all the elements of the set \verb!s! 
that are not in the set \verb!t!. These functions have the following signatures:
\begin{Verbatim}
def union(s: Set, t: Set): Set
def intersect(s: Set, t: Set): Set
def diff(s: Set, t: Set): Set	
\end{Verbatim}
\part
Define the function \verb!filter! which selects only the elements of a set that are accepted 
by a given predicate \verb!p!. The filtered elements are returned as a new set. 
The signature of \verb!filter! is as follows:
\begin{Verbatim}
def filter(s: Set, p: Int => Boolean): Set	
\end{Verbatim}
\end{parts}

\question
In this exercise, we are interested in functions used to make requests on elements of a set. 
The first function tests whether a given predicate is true for all elements of the set. 
This \verb!forall! function has the following signature:
\begin{Verbatim}
def forall(s: Set, p: Int => Boolean): Boolean	
\end{Verbatim}
Note that there is no direct way to find which elements are in a set. 
\verb!contains! only allows to know whether a given element is included. 
Thus, if we wish to do something to all elements of a set, then we have to iterate over all
integers, testing each time whether it is included in the set, and if so, 
to do something with it. Here, we consider that an integer \verb!x! 
has the property \verb!-1000 <= x <= 1000! to limit the search space.
\begin{parts}
\part
Implement \verb!forall! using tail-recursive (linear) recursion. 
For this, use a helper function nested in \verb!forall!. 
Its structure is as follows (you should replace the \verb!???! with your code):	
\begin{Verbatim}
def forall(s: Set, p: Int => Boolean): Boolean = {
 def iter(a: Int): Boolean = {
   if (???) ???
   else if (???) ???
   else iter(???)
 }
 iter(???)
}
\end{Verbatim}
\part
Using \verb!forall!, implement a function \verb!exists! which tests whether a set 
contains at least one element for which the given predicate is true. 
Note that the functions \verb!forall! and \verb!exists! behave like the 
\emph{universal} and \emph{existential} quantifiers of first-order logic.
\begin{Verbatim}
def exists(s: Set, p: Int => Boolean): Boolean	
\end{Verbatim}
\part
Finally, write a function \verb!map! which transforms a given set into another one by 
applying to each of its elements the given function. \verb!map! has the following signature:
\begin{Verbatim}
def map(s: Set, f: Int => Int): Set	
\end{Verbatim}
\end{parts}

\section*{Considerations}
\begin{itemize}
\item You will need to write anonymous functions in Scala to solve this exercise.
\item Sets are represented as functions. Think about what it means for an element 
to belong to a set, in terms of function evaluation. 
For example, how do you represent a set that contains all numbers between 1 and 100?
\item Most of the solutions for this exercise can be written as one-liners. If your code
is significantly longer then you probably need to rethink your solution. 
\item This exercise therefore requires more thinking (whiteboard, pen and paper) than coding ;-).
\item If you are having some trouble with terminology, have a look at the glossary at
	\url{http://docs.scala-lang.org/glossary/}.
\end{itemize}

\end{questions}

\end{document}